%=====================================================================
\chapter{Model Evaluation and Selection}\label{ch:evaluation}
%=====================================================================
\locator{3}{Part III --- Machine Learning Core}

\begin{outcomes}
By the end of this chapter you will be able to:
\begin{tight}
  \item \textbf{Build} a confusion matrix and derive accuracy, precision, recall, specificity and $F_1$ from it.
  \item \textbf{Explain} why accuracy is misleading under class imbalance, and choose an appropriate metric instead.
  \item \textbf{Read} an ROC curve and a precision--recall curve, and say which to use when.
  \item \textbf{Select} a decision threshold from the costs of the two error types.
  \item \textbf{Apply} k-fold cross-validation and explain why it beats a single split.
\end{tight}
\end{outcomes}

\begin{prereq}
\chref{probability} (the base-rate arithmetic --- precision \emph{is} the posterior
computed there), \chref{mlfundamentals} (splits) and \chref{classification}
(thresholds).
\end{prereq}

\begin{keyterms}
confusion matrix $\bullet$ true/false positive and negative $\bullet$ accuracy
$\bullet$ precision $\bullet$ recall (sensitivity) $\bullet$ specificity $\bullet$
$F_1$ score $\bullet$ ROC curve $\bullet$ AUC $\bullet$ precision--recall curve
$\bullet$ class imbalance $\bullet$ threshold tuning $\bullet$ k-fold
cross-validation $\bullet$ stratification
\end{keyterms}

\section{Why This Chapter Is the Hinge of Part III}

You cannot improve what you cannot measure, and you cannot choose between the
models of \chref{classification} without a number that reflects what you actually
care about. Every remaining chapter of this book assumes you can evaluate honestly.
Most published failures of applied machine learning are failures of evaluation, not
of algorithms.

\section{The Confusion Matrix}

\begin{definitionbox}[Confusion Matrix]
A table cross-tabulating predictions against truth. For two classes it has four
cells: \term{true positives} (TP), \term{false positives} (FP), \term{false
negatives} (FN) and \term{true negatives} (TN). Every classification metric is
some ratio of these four numbers.
\end{definitionbox}

\dsfig{%
\begin{tikzpicture}[font=\scriptsize]
\node[rounded corners=3pt, fill=primary, text=white, font=\scriptsize\bfseries,
      minimum width=5.6cm, minimum height=0.6cm] at (3.0,2.35) {ACTUAL TRUTH};
\node[font=\scriptsize\bfseries, text=primary] at (1.4,1.72) {Positive};
\node[font=\scriptsize\bfseries, text=primary] at (4.6,1.72) {Negative};
\node[rotate=90, font=\scriptsize\bfseries, text=white, fill=primary, rounded corners=3pt,
      minimum width=3.5cm, minimum height=0.6cm] at (-1.85,-0.15) {MODEL SAYS};
\node[font=\scriptsize\bfseries, text=primary, align=center, text width=1.3cm] at (-0.85,0.85) {Positive};
\node[font=\scriptsize\bfseries, text=primary, align=center, text width=1.3cm] at (-0.85,-1.15) {Negative};

\node[rounded corners=3pt, fill=greenhl!16, draw=greenhl, minimum width=3.0cm,
      minimum height=1.6cm, align=center] at (1.4,0.85)
     {{\color{greenhl!45!black}\bfseries TRUE POSITIVE}\\[2pt]{\tiny caught it, correctly}};
\node[rounded corners=3pt, fill=accent!13, draw=accent, minimum width=3.0cm,
      minimum height=1.6cm, align=center] at (4.6,0.85)
     {{\color{accent}\bfseries FALSE POSITIVE}\\[2pt]{\tiny false alarm}\\{\tiny (Type I)}};
\node[rounded corners=3pt, fill=accent!13, draw=accent, minimum width=3.0cm,
      minimum height=1.6cm, align=center] at (1.4,-1.15)
     {{\color{accent}\bfseries FALSE NEGATIVE}\\[2pt]{\tiny missed it}\\{\tiny (Type II)}};
\node[rounded corners=3pt, fill=greenhl!16, draw=greenhl, minimum width=3.0cm,
      minimum height=1.6cm, align=center] at (4.6,-1.15)
     {{\color{greenhl!45!black}\bfseries TRUE NEGATIVE}\\[2pt]{\tiny correctly ignored}};

% --- metric formulas, tied to the cells
\node[anchor=north west, align=left, text width=6.6cm] at (6.7,2.1) {%
  {\color{primary}\bfseries Accuracy} $=\dfrac{\text{TP}+\text{TN}}{\text{all}}$
  \quad{\tiny how often right overall}\\[6pt]
  {\color{greenhl!45!black}\bfseries Precision} $=\dfrac{\text{TP}}{\text{TP}+\text{FP}}$
  \quad{\tiny of those flagged, how many were real}\\[6pt]
  {\color{secondary}\bfseries Recall} $=\dfrac{\text{TP}}{\text{TP}+\text{FN}}$
  \quad{\tiny of the real ones, how many we caught}\\[6pt]
  {\color{goldhl!50!black}\bfseries Specificity} $=\dfrac{\text{TN}}{\text{TN}+\text{FP}}$
  \quad{\tiny of the negatives, how many left alone}\\[6pt]
  {\color{purplehl}\bfseries $F_1$} $=2\cdot\dfrac{\text{prec}\cdot\text{rec}}{\text{prec}+\text{rec}}$
  \quad{\tiny harmonic mean of the two}};

\node[align=left, text width=13.2cm, font=\tiny, text=gray!55!black] at (5.0,-2.8)
     {\textbf{Memory aid.} Precision reads \emph{down} the ``model says positive'' row:
      \emph{when I raise an alarm, am I right?} Recall reads \emph{across} the ``actually
      positive'' column: \emph{of everything I should have caught, how much did I?}};
\end{tikzpicture}}%
{The confusion matrix and the five metrics derived from it. Learn which cells each metric
uses; every question in this chapter reduces to that.}{confusion-matrix}

\begin{worked}[The Same Model, Five Metrics]
An intrusion detector is run on 10{,}000 sessions, of which 100 are genuine
attacks. It flags 190 sessions, of which 80 are real attacks.

\smallskip
\textbf{Fill in the four cells:}
\begin{tight}
  \item TP $= 80$ (flagged and real)
  \item FP $= 190 - 80 = 110$ (flagged, not real)
  \item FN $= 100 - 80 = 20$ (real, not flagged)
  \item TN $= 10{,}000 - 80 - 110 - 20 = 9{,}790$
\end{tight}

\smallskip
\textbf{Accuracy} $= \dfrac{80 + 9790}{10000} = \dfrac{9870}{10000} = \mathbf{98.7\%}$

\textbf{Precision} $= \dfrac{80}{80+110} = \dfrac{80}{190} = \mathbf{42.1\%}$

\textbf{Recall} $= \dfrac{80}{80+20} = \dfrac{80}{100} = \mathbf{80.0\%}$

\textbf{Specificity} $= \dfrac{9790}{9790+110} = \mathbf{98.9\%}$

\textbf{$F_1$} $= 2 \times \dfrac{0.421 \times 0.800}{0.421 + 0.800}
= 2 \times \dfrac{0.337}{1.221} = \mathbf{0.552}$

\smallskip
\textbf{The headline number is the useless one.} 98.7\% accuracy sounds superb ---
but a model that flags \emph{nothing at all} scores 99.0\% here, beating it.
Accuracy is dominated by the 9{,}900 easy negatives.

The informative pair is precision 42\% and recall 80\%: we catch four attacks in
five, and roughly three in five alarms are wasted analyst time. Whether that is
acceptable is the expected-value question of \chref{probability}, and it is the
only version of the result worth reporting.
\end{worked}

\begin{alertbox}[The Accuracy Paradox]
When one class is rare, always compute the majority-class baseline first. If 99\%
of traffic is benign, ``predict benign'' scores 99\%. Any accuracy figure below
that is worse than doing nothing, and any figure just above it may be worthless.
Report precision and recall, or a metric built from them.
\end{alertbox}

\section{Trading Precision Against Recall}

\begin{conceptbox}[They Move in Opposite Directions]
Lower the threshold: you flag more things, so recall rises and precision falls.
Raise it: precision rises and recall falls. You cannot maximise both, so you must
decide which error is more expensive --- and that is a domain decision, made with
the stakeholder.
\end{conceptbox}

\dsfig{%
\begin{tikzpicture}
\begin{groupplot}[
  group style={group size=2 by 1, horizontal sep=1.5cm},
  width=6.6cm, height=5.0cm,
  axis lines=left, tick label style={font=\tiny}, label style={font=\tiny},
  title style={font=\scriptsize\bfseries},
  legend style={font=\tiny, draw=gray!40},
]
\nextgroupplot[title={Metrics against threshold},
  xlabel={decision threshold}, ylabel={score},
  xmin=0, xmax=1, ymin=0, ymax=1.05,
  legend style={at={(0.5,-0.30)}, anchor=north, legend columns=3}]
\addplot[greenhl, line width=1.2pt, domain=0.03:0.97, samples=60] {0.06+0.88*x^0.75};
\addlegendentry{precision}
\addplot[secondary, line width=1.2pt, domain=0.03:0.97, samples=60] {1-0.93*x^1.7};
\addlegendentry{recall}
\addplot[purplehl, line width=1.2pt, dashed, domain=0.03:0.97, samples=60]
  {2*((0.06+0.88*x^0.75)*(1-0.93*x^1.7))/((0.06+0.88*x^0.75)+(1-0.93*x^1.7))};
\addlegendentry{$F_1$}
\draw[accent, dashed, line width=0.9pt] (axis cs:0.5,0) -- (axis cs:0.5,1.0);
\node[font=\tiny, accent, anchor=south, rotate=90] at (axis cs:0.5,0.06) {default 0.5};

\nextgroupplot[title={ROC curve},
  xlabel={false positive rate ($1-$specificity)}, ylabel={true positive rate (recall)},
  xmin=0, xmax=1, ymin=0, ymax=1.05,
  legend style={at={(0.99,0.16)}, anchor=south east}]
\addplot[gray!60, dashed, line width=0.9pt] coordinates {(0,0)(1,1)};
\addlegendentry{random (AUC 0.5)}
\addplot[greenhl, line width=1.4pt, domain=0:1, samples=90] {x^0.22};
\addlegendentry{good (AUC 0.90)}
\addplot[goldhl, line width=1.2pt, domain=0:1, samples=90] {x^0.52};
\addlegendentry{fair (AUC 0.70)}
\addplot[accent, only marks, mark=*, mark size=1.7pt] coordinates {(0.011,0.80)};
\node[font=\tiny, accent, anchor=west] at (axis cs:0.055,0.78) {worked example};
\node[font=\tiny, text=gray!55!black, anchor=north west] at (axis cs:0.03,0.99) {perfect $\rightarrow$};
\end{groupplot}
\end{tikzpicture}}%
{Left: every metric depends on the threshold, so quoting one without the threshold is
incomplete. Right: the ROC curve sweeps all thresholds at once; AUC is the area beneath
it.}{precision-recall-roc}

\begin{definitionbox}[ROC, AUC and the PR Curve]
The \term{ROC curve} plots recall against false-positive rate across all
thresholds. \term{AUC} is the area under it: the probability that the model ranks a
random positive above a random negative. 0.5 is chance, 1.0 is perfect.\\[3pt]
The \term{precision--recall curve} plots precision against recall instead.
\end{definitionbox}

\begin{alertbox}[Use PR, Not ROC, When Positives Are Rare]
ROC's $x$-axis is the false-positive \emph{rate} --- FP divided by the huge number
of negatives. With 9{,}900 negatives, 110 false positives is a rate of 0.011, which
looks negligible on an ROC plot even though it means 58\% of your alerts are wrong.
ROC curves flatter rare-event detectors badly. The PR curve, which uses precision
directly, does not. In cybersecurity, fraud and predictive maintenance, report the
PR curve.
\end{alertbox}

\section{Choosing the Threshold from Costs}

\begin{worked}[Setting a Threshold with Money]
A dropout-prediction model is to drive outreach calls.
\begin{tight}
  \item Cost of a false positive: one wasted 20-minute call $\approx$ \$15.
  \item Cost of a false negative: a student lost, worth $\approx$ \$4{,}000 in fees
        and outcomes.
\end{tight}
The ratio is roughly 1:267 --- a miss is 267 times more costly than a false alarm.

\smallskip
\textbf{Consequence.} You should accept a great many false alarms to avoid one
miss. Optimise recall subject to a staffing limit, not $F_1$ (which weights
precision and recall equally and would be the wrong objective here).

\textbf{Total expected cost at threshold $t$:}
\[
C(t) = 15 \times \text{FP}(t) + 4000 \times \text{FN}(t)
\]
Evaluate on the validation set across a grid of $t$ and pick the minimum.

\smallskip
\begin{center}
\begin{tabular}{@{}lrrrr@{}}
\toprule
\textbf{Threshold} & \textbf{FP} & \textbf{FN} & \textbf{Cost} & \\
\midrule
0.10 & 420 & \phantom{0}4 & \$22{,}300 & \\
0.20 & 210 & \phantom{0}9 & \$39{,}150 & \\
0.30 & 120 & 16 & \$65{,}800 & \\
0.50 & \phantom{0}48 & 31 & \$124{,}720 & \emph{the default} \\
0.05 & 780 & \phantom{0}2 & \$19{,}700 & $\leftarrow$ \textbf{minimum} \\
\bottomrule
\end{tabular}
\end{center}

\smallskip
\textbf{The optimum is 0.05}, not 0.5 --- a threshold that flags roughly a third of
the cohort. That is only workable if staff can handle the volume, which is why the
real constraint is usually written as ``maximise recall subject to flagging at most
20\% of students''. Either way, the default threshold was ten times too high and
cost six times more.
\end{worked}

\section{Cross-Validation}

\begin{definitionbox}[k-Fold Cross-Validation]
Split the training data into $k$ equal folds. Train on $k-1$ and validate on the
held-out one; repeat $k$ times so every fold is held out exactly once. Report the
mean and standard deviation of the $k$ scores.
\end{definitionbox}

\dsfig{%
\begin{tikzpicture}[font=\scriptsize]
\foreach \r in {1,...,5}{
  \pgfmathsetmacro{\yy}{-(\r-1)*0.72}
  \node[anchor=east, font=\tiny, text=gray!55!black] at (-0.25,\yy) {fold \r};
  \foreach \c in {1,...,5}{
    \pgfmathsetmacro{\xx}{(\c-1)*2.15}
    \ifnum\c=\r
      \node[rounded corners=2pt, fill=goldhl!30, draw=goldhl, minimum width=2.0cm,
            minimum height=0.56cm, font=\tiny\bfseries, text=goldhl!45!black] at (\xx,\yy) {VALIDATE};
    \else
      \node[rounded corners=2pt, fill=primary!12, draw=primary!50, minimum width=2.0cm,
            minimum height=0.56cm, font=\tiny, text=primary!55!black] at (\xx,\yy) {train};
    \fi}
  \node[anchor=west, font=\tiny, text=gray!55!black] at (9.3,\yy)
       {score$_\r$};}

\draw[decorate, decoration={brace, amplitude=4pt}, gray!60]
     (10.55,0.3) -- (10.55,-3.18)
     node[midway, right=5pt, align=left, font=\tiny, text=gray!55!black]
     {report \textbf{mean} $\pm$ \textbf{sd}\\[2pt]
      a large sd means the\\estimate is unstable ---\\usually too little data};

\node[anchor=north west, align=left, text width=10.5cm, font=\tiny, text=accent!75!black]
     at (-1.1,-3.75)
     {\textbf{Every fold must repeat the whole pipeline.} Scaling, imputation and feature
      selection are refitted inside each fold, on that fold's training part only. Fitting
      them once outside the loop leaks the validation data into training --- the error of
      \dsref{leakage-split}, hidden one level deeper.};
\end{tikzpicture}}%
{Five-fold cross-validation. Every observation is used for validation exactly once, so the
performance estimate uses all the data and comes with a measure of its own
stability.}{cross-validation}

\rowcolors{2}{white}{lightbg}
\begin{longtable}{@{}p{3.6cm}p{5.0cm}p{5.9cm}@{}}
\toprule
\rowcolor{primary!15}
\textbf{Variant} & \textbf{What it does} & \textbf{Use when} \\
\midrule
\endhead
k-fold ($k=5$ or $10$) & Standard random folds & The default for independent rows \\
Stratified k-fold & Preserves class proportions in every fold & Always, for classification --- especially imbalanced \\
Group k-fold & Keeps all rows of an entity in one fold & Multiple rows per student, device or user \\
Time-series split & Trains on the past, validates on the future & Any temporal data (\chref{timeseries}) \\
Leave-one-out & $k = n$ & Very small datasets; expensive and high-variance \\
\bottomrule
\end{longtable}

\begin{pitfall}
\begin{tight}
  \item \textbf{Reporting accuracy on imbalanced data.} Always compare to the
        majority baseline first.
  \item \textbf{Quoting a metric without its threshold.} ``Precision 0.9'' is
        meaningless alone.
  \item \textbf{Using ROC-AUC for rare positives.} Use PR-AUC.
  \item \textbf{Random cross-validation on time series or grouped rows.} Use
        \texttt{TimeSeriesSplit} or \texttt{GroupKFold}.
  \item \textbf{Tuning against the test set.} That is what the validation folds are
        for; the test set is used once (\chref{mlfundamentals}).
\end{tight}
\end{pitfall}

%---------------------------------------------------------------------
\begin{fourlenses}{Choosing a Metric}
\lensLA{\textbf{Optimise recall} at a fixed staffing budget. A missed at-risk
student is far costlier than an unnecessary supportive phone call, as the worked
example quantified. Report precision alongside it so staff know what proportion of
calls will find someone who was fine --- if that proportion is too embarrassing,
the programme will quietly stop making the calls regardless of what your metric
says.}

\lensPM{\textbf{Use $F_1$ or balanced accuracy}: over- and under-forecasting a
sprint carry roughly comparable costs, so neither error dominates. Beware tiny
validation folds --- with 40 sprints, 5-fold cross-validation validates on eight
sprints at a time, so report the standard deviation across folds and expect it to
be large.}

\lensIOT{\textbf{Optimise recall, then constrain precision by maintenance
capacity.} A missed failure stops a line; an unnecessary inspection costs an hour.
But note the second constraint: too many false alarms and technicians begin
ignoring the system, which converts a precision problem into a recall problem in
the real world. Use PR curves, since failures are rare.}

\lensSEC{\textbf{Use the PR curve and report precision at fixed recall}, e.g.
``precision 0.42 at 80\% recall''. ROC-AUC will look excellent and mean nothing,
for the reason in the alert box above. Also track \emph{alerts per analyst per
hour} --- a metric that appears in no textbook and decides whether the system is
used. And set the threshold from the cost ratio, not from 0.5, because in this
domain the two errors differ by orders of magnitude.}
\end{fourlenses}

%---------------------------------------------------------------------
\begin{lab}[Evaluate Honestly]
\begin{lstlisting}[language=Python]
import numpy as np
from sklearn.metrics import (confusion_matrix, classification_report,
                             roc_auc_score, average_precision_score,
                             precision_recall_curve)
from sklearn.model_selection import StratifiedKFold, cross_val_score

# --- 0. BASELINE FIRST -------------------------------------------------
print("majority-class accuracy:", max(np.mean(y_te), 1 - np.mean(y_te)))

# --- 1. the four cells -------------------------------------------------
proba = model.predict_proba(X_te)[:, 1]
pred  = (proba > 0.5).astype(int)
tn, fp, fn, tp = confusion_matrix(y_te, pred).ravel()
print(f"TP={tp}  FP={fp}  FN={fn}  TN={tn}")
print(classification_report(y_te, pred, digits=3))

# --- 2. threshold-free comparison; prefer PR-AUC when positives are rare
print("ROC-AUC:", roc_auc_score(y_te, proba))
print("PR-AUC :", average_precision_score(y_te, proba))   # <- the honest one

# --- 3. choose the threshold from COSTS, not from 0.5 ------------------
COST_FP, COST_FN = 15, 4000
best_t, best_cost = 0.5, float("inf")
for t in np.linspace(0.01, 0.95, 95):
    p = (proba > t).astype(int)
    _, fp_, fn_, _ = confusion_matrix(y_te, p).ravel()[[0,1,2,3]]
    cost = COST_FP*fp_ + COST_FN*fn_
    if cost < best_cost:
        best_t, best_cost = t, cost
print(f"cost-optimal threshold {best_t:.2f} (cost ${best_cost:,.0f})")

# --- 4. cross-validate, stratified, whole pipeline inside the folds ----
cv = StratifiedKFold(5, shuffle=True, random_state=0)
s = cross_val_score(pipeline, X_tr, y_tr, cv=cv, scoring="average_precision")
print(f"CV PR-AUC {s.mean():.3f} +/- {s.std():.3f}")
\end{lstlisting}
\end{lab}

%---------------------------------------------------------------------
\begin{checkpoint}
\begin{tightnum}
  \item A dataset is 97\% negative. A model scores 96.5\% accuracy. Is it any good?
  \item Of 500 flagged transactions, 60 were fraudulent; 20 frauds were missed.
        Compute precision and recall.
  \item You are told ``AUC is 0.94'' for a detector where 0.2\% of events are
        attacks. What do you ask for instead, and why?
\end{tightnum}
\end{checkpoint}

%---------------------------------------------------------------------
\begin{chaptersummary}
\begin{tight}
  \item Every classification metric is a ratio of the four confusion-matrix cells.
  \item Accuracy is dominated by the majority class; under imbalance it is
        actively misleading. Always compute the majority baseline first.
  \item Precision answers ``when I alarm, am I right?''; recall answers ``of the
        real cases, how many did I catch?'' They trade off against each other.
  \item ROC-AUC is threshold-free but flatters rare-event detectors; prefer PR-AUC
        when positives are rare.
  \item The 0.5 threshold is a default. Choose it by minimising expected cost, or
        by maximising recall subject to a capacity constraint.
  \item Cross-validate with stratification, group by entity, respect time order,
        and refit the entire pipeline inside each fold.
\end{tight}
\end{chaptersummary}

\begin{reviewq}
  \item \textbf{(MCQ)} Precision is: (a)~TP/(TP+FN) (b)~TP/(TP+FP)
        (c)~(TP+TN)/all (d)~TN/(TN+FP).
  \item \textbf{(MCQ)} When positives make up 0.5\% of the data, the most
        informative curve is: (a)~ROC (b)~precision--recall (c)~learning curve
        (d)~residual plot.
  \item \textbf{(Short)} Explain the accuracy paradox with a numeric example.
  \item \textbf{(Short)} Why must feature scaling be refitted inside every
        cross-validation fold?
  \item \textbf{(Short)} Give a scenario where you would deliberately accept a
        precision of 0.2, and justify it.
  \item \textbf{(Applied)} A model tested on 20{,}000 records with 400 positives
        flags 900 records, 300 of which are true positives. Build the full confusion
        matrix and compute accuracy, precision, recall, specificity and $F_1$.
        Then state which two numbers you would put in an executive summary and why.
  \item \textbf{(Applied)} For the same model, a false positive costs \$8 and a
        false negative costs \$1{,}200. Explain how you would find the optimal
        threshold, and predict whether it will be above or below 0.5, with reasoning.
\end{reviewq}
